\chapter{The Rectangle Method}

\section{Rectangles and Partitions}

\begin{definition}[Rectangle]
  \label{def:rectangle}
  \lean{DetProtocol.isRectangle}
  \leanok
  A set $R \subseteq X \times Y$ is a \emph{combinatorial rectangle} if
  $R = A \times B$ for some $A \subseteq X$ and $B \subseteq Y$.
  Equivalently, whenever $(x,y),(x',y') \in R$ we also have $(x',y),(x,y') \in R$.
\end{definition}

\begin{lemma}[Rectangle characterization]
  \label{lem:rectangle_iff}
  \lean{DetProtocol.isRectangle_iff}
  \leanok
  \uses{def:rectangle}
  $R$ is a rectangle if and only if for all $(x,y),(x',y') \in R$,
  the mixed pairs $(x',y)$ and $(x,y')$ also belong to $R$.
\end{lemma}

\begin{definition}[Monochromatic rectangle partition]
  \label{def:mono_rect_partition}
  \lean{DetProtocol.isMonoRectPartition}
  \leanok
  \uses{def:rectangle}
  A collection $\mathcal{P}$ of subsets of $X \times Y$ is a
  \emph{monochromatic rectangle partition} for $g : X \to Y \to \alpha$ if:
  every $R \in \mathcal{P}$ is a rectangle, every $R \in \mathcal{P}$ is
  $g$-monochromatic (all pairs in $R$ have the same $g$-value), $\mathcal{P}$
  covers $X \times Y$, and distinct members of $\mathcal{P}$ are disjoint.
\end{definition}

\begin{definition}[Leaf rectangles]
  \label{def:leaf_rectangles}
  \lean{DetProtocol.leafRectangles}
  \leanok
  \uses{def:det_protocol, def:rectangle}
  The \emph{leaf rectangles} of a protocol $\pi$ are the sets of inputs that
  reach each leaf of the protocol tree.  Each leaf rectangle is a combinatorial
  rectangle obtained by intersecting the half-spaces cut by each bit sent along
  the path to that leaf.
\end{definition}

\begin{lemma}[Lemma 1.4: leaf rectangles partition the input space]
  \label{lem:leaf_partition}
  \lean{DetProtocol.leafRectangles_cover}
  \lean{DetProtocol.leafRectangles_disjoint}
  \lean{DetProtocol.leafRectangles_isRectangle}
  \leanok
  \uses{def:leaf_rectangles, def:mono_rect_partition}
  For any protocol $\pi$, the leaf rectangles of $\pi$ form a partition of
  $X \times Y$ into combinatorial rectangles.
  Moreover, if $\pi$ computes $g$, each leaf rectangle is $g$-monochromatic.
\end{lemma}

\begin{lemma}
  \label{lem:leaf_card}
  \lean{DetProtocol.leafRectangles_card}
  \leanok
  \uses{def:leaf_rectangles, def:protocol_run}
  A protocol of complexity $c$ has at most $2^c$ leaf rectangles.
\end{lemma}

\begin{theorem}[Theorem 1.6: rectangle partition bound]
  \label{thm:rect_partition}
  \lean{DetProtocol.rectangle_partition}
  \lean{DetProtocol.mono_rectangle_partition_of_det_cc}
  \leanok
  \uses{lem:leaf_partition, lem:leaf_card, def:mono_rect_partition}
  If $\pi$ computes $g$ and has complexity $c$, then the input space $X \times Y$
  can be partitioned into at most $2^c$ monochromatic rectangles with respect to $g$.
  Equivalently, if $\Det(g) \le n$ then there is a monochromatic rectangle
  partition of size at most $2^n$.
\end{theorem}

\begin{theorem}[Rectangle lower bound method]
  \label{thm:rect_lower_bound}
  \lean{DetProtocol.det_cc_lower_bound}
  \leanok
  \uses{thm:rect_partition, def:det_cc}
  To show $\Det(g) \ge n+1$, it suffices to show that every monochromatic
  rectangle partition of $g$ has more than $2^n$ members.
\end{theorem}
